\section{Related Work}
\label{sec:related}

\paragraph{IPI attacks and benchmarks.}
Indirect prompt injection was introduced by \citet{greshake2023signed}; \citet{liu2023formalizing} formalize
prompt-injection attacks and defenses. BIPIA \citep{yi2023bipia} inserts attack strings into email, table, code, web
and summarization contexts at the start, middle or end of the document. InjecAgent
\citep{zhan-etal-2024-injecagent} places attacker instructions in the tool responses of tool-using agents, and
AgentDojo \citep{debenedetti2024agentdojo} and $\tau$-bench \citep{yao2024taubench} provide agent environments whose
tool outputs can carry injections. Our artifact audit concerns how BIPIA-style insertion data is
\emph{constructed}, not the attacks themselves.

\paragraph{Detectors.}
Content-only classifiers include PIGuard \citep{li-etal-2025-piguard}, which targets over-defense on trigger words,
Sentinel \citep{ivry2025sentinel}, and Prompt Guard 2 in LlamaFirewall \citep{chennabasappa2025llamafirewall}.
DataSentinel \citep{liu2025datasentinel} detects injections with a fine-tuned LLM through a known-answer test, and
Attention Tracker \citep{hung2024attentiontracker} scores attention distraction inside the model. PromptArmor
\citep{shi2025promptarmor} prompts an off-the-shelf LLM to detect and remove injected prompts. MELON
\citep{zhu2025melon} compares an agent's actions with and without the user task; it defends agent trajectories
rather than classifying text.

\paragraph{How detectors are evaluated.}
\citet{liu2026passing} re-evaluates public detectors and two task-aware LLM judges on AgentDojo and $\tau$-bench tool
outputs replayed without an LLM. Detector rankings on BIPIA do not carry over to the agent benchmarks, the best BIPIA
detector was trained on BIPIA attack strings, and the form of a detector's training inputs, more than attack-string
overlap, determines where it works. \citet{fomin2026benchmarks} shows that in-distribution splits overstate the
generalization of malicious-prompt classifiers and proposes leave-one-dataset-out evaluation, and
\citet{akinrele2026regime} find that detector rankings depend on the deployment regime and threshold. We build on
\citet{liu2026passing} directly: we rebuild its agent-benchmark sets with its released scripts, and ask a question it
leaves open, namely whether removing construction artifacts from the training data changes transfer.

\paragraph{Shortcuts in IPI detection.}
\citet{li2026surface} show that defenses against prompt attacks learn surface heuristics, including position,
trigger-token and topic biases, and name counterfactual benign examples as future work. \citet{kale2026probes} show
that on AgentDojo, where an injection sits in a conversation and how long the conversation is nearly solve the
standard probe evaluation, and that training with matched negatives (the same conversation with the injection
removed, and harmless look-alikes) reduces false alarms. Our artifact analysis is confirmatory for the general
claim that IPI evaluations contain content-free shortcuts. It is new in measuring a specific inventory for
BIPIA-style insertion data and in rebuilding the data under pre-registered gates.

\paragraph{Task-, intent- and alignment-aware detection.}
Several methods ask whether an instruction in external data deviates from the user's task. TaskTracker
\citep{abdelnabi2024drift} detects task drift from activation deltas, PromptSleuth \citep{wang2025promptsleuth}
reasons over task-level intent, and IntentGuard \citep{kang2025intentguard} analyzes a reasoning model's
instruction-following intent. The closest work is AlignSentinel \citep{jia2026alignsentinel}, a three-class
classifier (misaligned instruction, aligned instruction, no instruction) over a backend LLM's attention maps,
trained on a GPT-4o-generated benchmark. In that benchmark, aligned and misaligned inputs are different texts under
the same user prompt and tool; we found no setting in which identical instruction text is aligned under one request
and misaligned under another, and no released code, data or model. \citet{liu2026passing} also gives a small
classifier the user goal as an extra segment and finds that it does not help; that classifier was never trained to
use the goal. Conditioning a detector on the request is therefore not new. What is new here is the evaluation
design that varies only the request, and its use as a training signal for a black-box text classifier.

\paragraph{Counterfactual data and artifacts.}
Our pairs apply contrast sets \citep{gardner-etal-2020-evaluating} and counterfactually augmented data
\citep{kaushik2019counterfactual} to IPI detection. The artifact audit follows work on annotation artifacts
\citep{gururangan-etal-2018-annotation}, syntactic heuristics \citep{mccoy-etal-2019-right} and shortcut learning
in general \citep{geirhos2020shortcut}.

\paragraph{Which instructions in data are injections.}
A strict view treats every instruction in external data as untrusted: StruQ \citep{chen2024struq} trains a model to
follow only instructions in the prompt portion of its input. The instruction hierarchy \citep{wallace2024hierarchy}
instead defines behavior for instructions of different privilege that \emph{conflict}, and teaches models to ignore
lower-privileged instructions selectively. Our definition (\S\ref{sec:task}) is closer to the latter.
